\clearpage
\chapter*{ЗАКЛЮЧЕНИЕ}
\addcontentsline{toc}{chapter}{ЗАКЛЮЧЕНИЕ}

В ходе выполнения выпускной квалификационной работы разработаны вычислительная модель и программный комплекс для расчёта термодинамических свойств воды и водяного пара по стандарту IAPWS-IF97 \cite{iapws_if97_2012}.  
Достигнута поставленная цель --- создано единое вычислительное ядро и набор каналов доставки результатов, ориентированных как на интерактивные инженерные расчёты, так и на интеграцию в сервисную инфраструктуру.  

В рамках работы решены основные задачи, сформулированные во введении.  
Реализовано вычислительное ядро \texttt{if97\_core} на языке Rust, включающее региональные уравнения IF97, маршрутизацию расчётов, централизованную валидацию входных данных и унифицированный формат результата \cite{rust_book}.  
Поддержаны прямые и обратные маршруты расчёта по основным входным парам параметров ($pT$, $pH$, $pS$, $pX$) и расширение по паре $(\rho, T)$, реализованное как численный подбор давления методом секущих с последующим расчётом по стандартному маршруту $pT$.  
Для обеспечения воспроизводимости и производительности коэффициенты стандарта организованы как набор CSV-таблиц с последующей build-time кодогенерацией статических массивов, что исключает runtime I/O и снижает риск ошибок ручного ввода.  

Разработаны две пользовательские оболочки, демонстрирующие переносимость архитектуры и практическую применимость результата.  
Нативное настольное приложение на FLTK обеспечивает одиночный и табличный режимы расчёта, построение термодинамических диаграмм и экспорт данных, что соответствует типовым инженерным сценариям \cite{fltk_rs_docs}.  
Desktop-оболочка на базе Yew/WASM+Tauri реализует разделение frontend и backend, использование общего слоя \texttt{if97\_app\_api} и интеграцию с функциями ОС (файловые диалоги, экспорт, журналы), сохраняя единый источник вычислений в \texttt{if97\_core} \cite{tauri_docs}.  

Сервисный контур реализован в виде gRPC-сервиса пакетных расчётов.  
Выбран потоковый сценарий взаимодействия и формат данных Struct-of-Arrays, упрощающий высокопроизводительную обработку батчей \cite{grpc_docs}.  
В сервисе применено разделение I/O и CPU-нагрузки (tokio и rayon), возврат результатов через \texttt{mpsc}-очередь и механизмы backpressure с ограничениями in-flight батчей.  
Эксплуатационная устойчивость обеспечена health-check через \texttt{tonic-health} и корректным завершением в режиме drain.  
Для эксплуатации подготовлены контейнеризация в Docker и манифесты Kubernetes, а также CI/CD на GitHub Actions, обеспечивающий воспроизводимую сборку и публикацию артефактов \cite{docker_docs,kubernetes_docs,github_actions_docs}.  
По состоянию отчёта \texttt{TESTS.md}
от 09.04.2026
все функциональные тесты проекта пройдены успешно,
а perf/stress-сценарии вынесены в отдельные Bench-секции.

Нагрузочные измерения подтвердили практическую применимость сервисного контура:
в серии из 6 запусков \texttt{grpc\_loadtest}
в \texttt{profile=debug}
(500 батчей по 2048 точек)
средний throughput составил $129632.83$ points/s
с 95\% доверительным интервалом
$[123299.91;\ 135965.75]$.
В длинном прогоне (\texttt{batches=2000})
получено $122865$ points/s.  

Практическая значимость выполненной работы заключается в том, что единая физическая модель и единый программный API доступны в нескольких формах поставки: локальный инженерный калькулятор, кроссплатформенная desktop-оболочка и сервис для пакетных расчётов.  
Архитектура «одно ядро --- несколько адаптеров» обеспечивает совпадение результатов во всех каналах доставки и снижает риск расхождений при дальнейшей доработке системы.  

В качестве направлений дальнейшего развития можно выделить:
\begin{itemize}
\item расширение набора поддерживаемых входных пар параметров и специализированных режимов расчёта для прикладных сценариев;  
\item углубление верификации ядра за счёт расширения набора контрольных точек и автоматизированных регрессионных тестов;  
\item развитие графического функционала (дополнительные диаграммы, изолинии, улучшение инструментов анализа и экспорта);  
\item развитие сервисного слоя (расширение API, дополнительные типы батчевых расчётов, эксплуатационные метрики и наблюдаемость).  
\end{itemize}
